\section{De las matemáticas al descubrimiento científico en sentido amplio}

AI for Math también es importante porque las matemáticas son la parte más formalizada y verificable del descubrimiento científico. Si la IA logra cerrar en matemáticas el ciclo «formular conjeturas--construir demostraciones--verificación por máquina--evaluación humana», métodos similares podrían trasladarse a la física, los materiales, la verificación de programas y el descubrimiento de algoritmos.

\subsection{Por qué las matemáticas son un campo de pruebas}

Las matemáticas tienen tres ventajas. Primera: la corrección puede verificarse de forma rigurosa mediante sistemas formales. Segunda: los datos históricos y las bibliotecas de teoremas se han acumulado de manera estructurada. Tercera: la comunidad matemática tiene una tradición de evaluación de la elegancia, el poder explicativo y la importancia que es clara, pero difícil de cuantificar. Esto hace que sean adecuadas para el entrenamiento de máquinas y, a la vez, que pongan en evidencia las insuficiencias de ese entrenamiento.

\begin{knowledgebox}{Tres tipos de retroalimentación en el descubrimiento científico}
Las matemáticas aportan verificación formal; el código aporta pruebas y retroalimentación de ejecución; las ciencias experimentales aportan retroalimentación de observaciones y experimentos. Los tres son entornos esenciales para que la IA pase de generar lenguaje a hacer descubrimientos reales. El Agent de EP139, el laboratorio de modelos de EP138 y AI for Math de EP137 buscan, en el fondo, entornos que proporcionen retroalimentación.
\end{knowledgebox}

\subsection{El destino final de la colaboración entre humanos y máquinas}

El destino ideal de AI for Math no es necesariamente «sustituir a los matemáticos». Una forma más realista y también más valiosa es que la IA ayude a los matemáticos a formalizar demostraciones engorrosas, buscar lemma, encontrar contraejemplos, organizar la bibliografía y proponer conjeturas candidatas; los matemáticos, por su parte, se encargan del criterio para elegir problemas, la elección de direcciones, la interpretación y el juicio final. Así, la IA se convierte en una compañera de investigación y no en un simple sistema que responde preguntas de forma automática.

\begin{importantbox}{De herramienta a compañera}
Cuando la IA solo realiza pasos locales, es una herramienta; cuando puede comprender el contexto, acumular experiencia, proponer direcciones candidatas y someterse a verificación, empieza a parecerse a una compañera de investigación. AI for Math es un modelo de alta dificultad de esta transformación.
\end{importantbox}

\subsection{Resumen de la sección}

Las matemáticas son un campo de pruebas de alto valor para la entrada de la IA en el descubrimiento científico. Su fuerte verificabilidad hace que el entrenamiento sea más claro, y su estética e intuición nos recuerdan que el descubrimiento científico no puede reducirse a un único indicador.
